% Figure: first three mode shapes of a string fixed at both ends.
% y_n(x) = sin(n pi x / L), L = 1, n = 1,2,3.
\begin{figure}[htbp]
\centering
\begin{tikzpicture}
\begin{axis}[
  width=12cm, height=6cm,
  xlabel={$x/L$}, ylabel={mode shape $\phi_n(x)$},
  xmin=0, xmax=1, ymin=-1.25, ymax=1.25,
  axis lines=left,
  domain=0:1, samples=300,
  legend pos=north east,
  legend cell align=left,
  every axis x label/.style={at={(current axis.right of origin)}, anchor=west},
  every axis y label/.style={at={(current axis.above origin)}, anchor=south},
  grid=major, grid style={enggray!20},
  tick label style={font=\scriptsize},
  label style={font=\small},
  legend style={font=\scriptsize},
]
\addplot[engblue, very thick] {sin(deg(pi*x))};
\addlegendentry{$n=1$ (fundamental)}
\addplot[engred, very thick] {sin(deg(2*pi*x))};
\addlegendentry{$n=2$ (one node)}
\addplot[engamber, very thick] {sin(deg(3*pi*x))};
\addlegendentry{$n=3$ (two nodes)}
\addplot[enggray, thin, forget plot] {0};
% mark interior nodes of n=2 and n=3
\node[engred, font=\scriptsize] at (axis cs:0.5,-0.12) {$\bullet$};
\node[engamber, font=\scriptsize] at (axis cs:0.3333,0.12) {$\bullet$};
\node[engamber, font=\scriptsize] at (axis cs:0.6667,-0.12) {$\bullet$};
\end{axis}
\end{tikzpicture}
\caption[Mode shapes of a fixed-fixed string]{The first three mode
shapes of a uniform string fixed at both ends,
$\phi_n(x)=\sin(n\pi x/L)$. The fundamental ($n=1$, blue) has no
interior node; the second mode ($n=2$, red) has one interior node at
the midpoint; the third ($n=3$, amber) has two interior nodes. A
point that stays still in a given mode is a node; a point of maximum
motion is an antinode. The natural frequencies follow the integer
ratio $1:2:3$, the harmonic series that gives a plucked string its
pitched, musical tone. Driving the string at one of its nodes for a
given mode fails to excite that mode at all, which is why a guitarist
who lightly touches the string at its midpoint silences the
fundamental and brings out the octave harmonic. The beam, by
contrast, has the same node-counting pattern but a non-harmonic
frequency ratio $1:4:9$. Computed from $\sin(n\pi x/L)$ with $L=1$.}
\label{fig:vol03ch06:mode-shapes}
\end{figure}
